\scanpage{94}

\[
\begin{aligned}
X_i\subset X_j&\Longleftrightarrow\forall i'\in\bar i,\quad X_{i'}^*\subset X_j=\bigcup_{i'\in\bar j}X_{i'}^*\\
&\Longleftrightarrow\forall i'\in\bar i,\text{ t.q. }X_{i'}^*\neq0,\quad\unsure{\text{on a}|\mathord{\cdot}\mathord{-}\mathord{\cdot}}\quad i'\in\bar j\text{ i.e. }i'\leq j
\end{aligned}
\]
Si la stratification est fidèle, alors on aura \unsure{donc|même}
\[
\begin{aligned}
(92\unsure{)}\qquad X_i\subset X_j\Longleftrightarrow X_i^*\subset X_j\Longleftrightarrow i\leq j&\\
(\text{si strat. fidèle})&
\end{aligned}
\]
ce qui implique que $I\longrightarrow\mathcal P(X)$ est injective, et que la relation d'ordre de $I$ est induite par la relation d'inclusion de $\mathcal P(X)$. On en peut identifier $I$ à une partie de $\mathcal P(X)$, et la donnée de la str. fidèle à la donnée d'un \deleted{\unsure{ens.|fam.}} \inserted{\unsure{rec.|rec} loc \unsure{fini|finie}} de parties fermées \inserted{non vides} de $X$, \inserted{les ``strates fermées'', de réunion $X$} satisfaisant la condition (Strat 3) \inserted{sur l'intersection,} et la condition qu'aucune strate fermée n'est réunion de ses sous-strates strictement plus petites. Si on \unsure{veut|est} Strat 5, cela s'exprime par l'axiome plus fort que dans tt strate fermée.

Si $J$ est une partie de $I$, on pose
\[
(93)\qquad\left\{
\begin{aligned}
X_J^*&=\bigcup_{j\in J}X_j^*\qquad(\text{pas néc fermé})\\
X_J&=\bigcup_{j\in J}X_j\qquad\text{fermé de }X
\end{aligned}\right.
\]
\margin{\fbox{\parbox{0.88\linewidth}{NB $J\unsure{\mapsto|\mathrel{\mapstochar\relbar}}X_J^*$ est compatible aux réunions et intersections $X_J^*\cap X_K^*=X_{J\cap K}^*$, $X_{J\cup K}^*=X_J^*$}}}
\deleted{\unsure{Nous|On} \unsure{utilisera|utiliserons} cette \unsure{notation|notations} \unsure{surtout} si $J$ est une bande. Alors on \unsure{aura|va}}
\margin{93}

